\chapter{数学语言、证明、反例与量纲}
\label{ch:language-proof-dimension}

量化研究里最危险的错误，往往不是把一个加法算错，而是把一句听起来合理的话当成了已经定义清楚的结论。“分散化降低风险”“相关性高意味着存在长期关系”“收益越高越好”都可以成为有用的研究假设，却还不是数学命题。它们没有说明对象、允许的范围、量词、风险口径和时间尺度，因此既不能严格证明，也不能设计有意义的反例。

本章建立一套之后每章都会用到、但不再重复讲解的语言。读完以后，读者应当能把研究直觉改写成可判定的命题，判断一个条件是必要还是充分，选择直接证明、逆否、反证或反例，并在任何计算开始前检查单位、定义域和时间尺度。后文的计算只重复已经由纸笔得到的有限结果；它不能替代证明。

\section{把研究直觉写成可证明命题}

\textbf{定义}规定一个术语在当前语境中的含义；\textbf{命题}是具有确定真假值的陈述。定义本身不是由样本“验证”的，命题则需要证明、反例，或者在无法得到解析结论时提供带误差声明的计算证据。

例如，组合简单收益
\[
  R_p=\sum_{i=1}^n w_i r_i
\]
是定义；“若 $w_i\geq0$ 且 $\sum_iw_i=1$，则 $R_p$ 不超过最大分量收益”才是命题。后一句已经指明了收益属于同一期间、权重在期初确定、结论对哪些收益向量成立以及“不超过”的精确定义。

把自然语言主张写成
\[
  \forall x\in\mathcal D,\qquad A(x)\Longrightarrow B(x)
\]
通常是最有用的第一步。$\mathcal D$ 是允许输入的集合，$A$ 汇总假设，$B$ 是结论。若主张只声称存在一个成功对象，则写成
\[
  \exists x\in\mathcal D\quad\text{使得}\quad B(x).
\]
这两个量词要求完全不同的证据：找到一个盈利参数足以支持存在性，却不能支持“所有参数都盈利”；在一百万个随机样本里没有发现失败，也不能证明全称命题。

把“价格上涨时策略一定盈利”写完整，可以得到：“对所有允许的价格路径，只要期末价格高于期初价格，扣除全部交易成本后的策略 P\&L 就为正。”现在反驳的任务变得清楚：构造一条价格上涨、但换手成本超过毛收益的合法路径即可。反例必须保留原命题的全部假设，只让结论失败；如果反例同时改变了价格口径和交易成本，就无法知道究竟是哪条条件造成了失败。

全称命题的否定也决定了反例的形状：
\[
  \neg\bigl[\forall x\in\mathcal D,\ A(x)\Rightarrow B(x)\bigr]
  \quad\Longleftrightarrow\quad
  \exists x\in\mathcal D:\ A(x)\text{ 成立且 }B(x)\text{ 不成立}.
\]
相反，否定“存在一个 $x$ 使 $B(x)$ 成立”需要证明对所有允许对象 $B$ 都失败，这通常比找到一个成功例子困难得多。

\section{必要、充分与证明的选择}

若 $A$ 推出 $B$，写作 $A\Rightarrow B$。此时 $A$ 是 $B$ 的充分条件，$B$ 是 $A$ 的必要条件；只有两个方向都成立，才可以写 $A\Leftrightarrow B$。可以把箭头翻译成研究语言：手续费检查可能是“回测可信”的必要条件，因为完全不计成本的报告必然不完整；但它不是充分条件，因为未来函数、样本选择或容量问题仍可能存在。

\begin{center}
\begin{tabular}{lll}
\toprule
表达 & 真正要证明的方向 & 反驳方法 \\
\midrule
$A$ 对 $B$ 充分 & $A\Rightarrow B$ & 找到 $A$ 真、$B$ 假 \\
$A$ 对 $B$ 必要 & $B\Rightarrow A$ & 找到 $B$ 真、$A$ 假 \\
$A$ 与 $B$ 充要 & 两个方向 & 破坏任一方向 \\
\bottomrule
\end{tabular}
\end{center}

一个容易混淆的例子是“矩阵正定所以可逆”。这条箭头成立，但反向不成立。二维旋转矩阵
\[
  Q=\begin{pmatrix}0&-1\\1&0\end{pmatrix}
\]
满足 $\det(Q)=1$，因而可逆；然而对任意非零 $x$ 都有 $x^\top Qx=0$，所以它不是正定矩阵。后面学习协方差矩阵和优化时，若把“可逆”误读成“正定”，就会错误地声称目标函数有唯一最小值。

证明方法应该由命题结构决定，而不是由篇幅决定。直接证明从假设出发逐步到达结论，适合不等式、集合包含和代数恒等式；逆否证明把 $A\Rightarrow B$ 改写成 $\neg B\Rightarrow\neg A$，适合结论失败更容易描述的场景；反证法假设结论不成立并推出明确矛盾，常用于唯一性和不存在性。每一步都要能回答“用了哪条假设”；如果一条假设从未出现，它要么多余，要么证明漏了一步。

反例和数值实验的职责不同。反例可以推翻全称命题，却不能证明另一个全称命题；有限实验可以检查固定算例、搜索反例或核对一个实现，但不能覆盖连续定义域。判断一个实验是否有意义，先问它的输入是否在命题定义域内，参照值是否由独立方法得到，失败是否只改变了待检验的假设。

% Generated by tools/build_figures.py; edit figures/figure-specs.json instead.
\MFQFigure{tex/figures/generated/upper-ch01-proof-logic.pdf}{证明链与反例链回答不同问题：证明从假设推出结论，反例则精确切断一个过强命题。}{upper-ch01-proof-logic}


\section{凸组合收益：一个完整的证明、几何和反例}

设 $r_i\in\mathbb R$ 是第 $i$ 个资产在同一期间、同一口径下的简单收益，$w_i$ 是期初权重。令
\[
  R_p=\sum_{i=1}^n w_i r_i,\qquad
  m=\min_i r_i,\qquad M=\max_i r_i.
\]
我们要证明的不只是“加权平均通常落在范围内”，而是一个带有精确必要条件的命题。

\begin{theorem}[凸组合界]
在 $\sum_iw_i=1$ 的前提下，以下两件事等价：
\begin{enumerate}
  \item 对所有实数向量 $(r_1,\ldots,r_n)$，都有 $m\leq R_p\leq M$；
  \item 对每个 $i$，都有 $w_i\geq0$。
\end{enumerate}
\end{theorem}

先证充分性。若 $w_i\geq0$，则每个 $r_i-m\geq0$，而
\[
 R_p-m
 =\sum_iw_ir_i-m\sum_iw_i
 =\sum_iw_i(r_i-m)\geq0.
\]
这里把 $m$ 写成 $m\sum_iw_i$ 的一步，正是归一化假设发挥作用的地方。同理，
\[
 M-R_p
 =M\sum_iw_i-\sum_iw_ir_i
 =\sum_iw_i(M-r_i)\geq0.
\]
于是 $m\leq R_p\leq M$。

再证必要性。假设某个 $w_j<0$。利用“对所有收益向量”这一量词，选择 $r_j=0$、其余 $r_i=1$。此时 $m=0$、$M=1$，且
\[
  R_p=\sum_{i\ne j}w_i=1-w_j>1=M.
\]
所以界失败。这个构造保留了权重和为 1，只删除了非负性，因而确实说明非负性是必要条件。

等号条件也可以从证明中读出：在非负权重下，$R_p=m$ 当且仅当每个正权重资产都取得 $m$；$R_p=M$ 的条件类似。只要两个正权重资产的收益不同，组合收益就严格位于区间内部。几何上，一维的所有凸组合构成区间；高维中同样的权重构成收益向量的凸包。负权重并非“错误”，但它允许越出凸包，因此需要另行声明杠杆、融资和风险约束。

现在代入一个可复核的数字例子：
\[
  r=(0.02,-0.01,0.03),\qquad w=(0.5,0.3,0.2).
\]
三个贡献分别为 $0.010,-0.003,0.006$，所以 $R_p=0.013$，落在 $[-0.01,0.03]$ 内。只删除非负性，改用
\[
  w'=(1.5,-0.5,0),\qquad \sum_iw'_i=1,
\]
则 $R'_p=0.035>0.03$。若使用 $w''=(2,0,0)$，虽然也会越界，但它同时破坏了权重和为 1，不能作为“非负性必要”的隔离证据。

\section{量纲、时间尺度与研究计算}

量纲是计算前最便宜、也最容易被忽略的检查。设持仓 $q=200\ \mathrm{share}$，价格变动 $\Delta p=0.5\ \mathrm{USD/share}$，则
\[
 [q\Delta p]
 =\mathrm{share}\times\frac{\mathrm{USD}}{\mathrm{share}}
 =\mathrm{USD},\qquad \mathrm{P\&L}=100\ \mathrm{USD}.
\]
若初始价格为 $p_0=100\ \mathrm{USD/share}$，资本基数为 $qp_0=20000\ \mathrm{USD}$，对应的简单收益才是无量纲的
\[
  \frac{q\Delta p}{qp_0}=\frac{\Delta p}{p_0}=0.005=0.5\%.
\]
把 $100$ 美元 P\&L 与 $2\%$ 收益相加，即使程序能够运行，也没有数学意义。

“无量纲”也不意味着任意可比。日收益和月收益没有货币单位，却属于不同持有期；小数 $0.02$ 与百分数 $2$ 相差百倍；对数收益与简单收益在大幅波动时也不能无条件互换。若在独立同分布、方差有限且口径稳定的强假设下年化日均值和波动率，缩放分别为
\[
  \mu_{\mathrm{ann}}\approx252\mu_{\mathrm{day}},\qquad
  \sigma_{\mathrm{ann}}\approx\sqrt{252}\,\sigma_{\mathrm{day}}.
\]
把均值乘 $\sqrt{252}$ 或把波动率乘 $252$，都是量纲错误；自相关、波动聚集和结构变化还会使形式正确的缩放公式失效。

本章的三组数字——组合收益 $0.013$、区间 $[-0.01,0.03]$ 和反例 $0.035$——应当先由纸笔得到，再由透明程序重复。程序可以检查向量长度、权重和、有限性、同一资产顺序和同一持有期，但数组运算不知道元素究竟是收益、美元金额还是错位日期的数据。独立参照值必须来自手算、解析式、枚举或分数运算，而不能由被测函数自己生成。

读者可以进一步改变一个条件，观察结论如何变化：让权重和不为 1，改变非负性，或把收益和权重长度改成不同值。每次只改一件事，才能把失败归因到具体假设。数值实验的正确结论是“这个实现通过了这些输入的核对”，而不是“程序证明了定理”。

\section{把判断写成可迁移的研究语言}

遇到一个新的研究主张时，可以按以下顺序思考，但不必把它写成每章重复的检查表。先明确对象和单位，再写量词和假设；随后选择证明或反例；最后才决定是否需要计算。这个顺序能把“策略有正样本外收益”拆成至少几个不同命题：收益是否定义在固定信息集上，是否扣除成本，是否在独立样本上评价，是否足以支持部署。通过其中一个必要检查，只能说明该检查没有否决结论，不能把它升级成充分证据。

后续章节会把同一语言分别用于极限、积分、矩阵、概率、统计和优化。读者不需要记住一套工程术语，只需要在每次使用定理前回答三个问题：对象属于哪个空间，假设在哪里被使用，结论比计算结果多说了什么。量纲一致、程序运行成功和样本中没有明显异常，都是必要时有用的局部证据，却都不是数学结论本身。

\section*{章末练习}
\begingroup\small
\subsection*{口述概念}
\begin{enumerate}[nosep,leftmargin=*]
  \item 区分定义与命题。为什么“组合收益定义为 $R_p=\sum_iw_ir_i$”不能用样本数据证明，而凸组合界需要证明？
  \item 在 $A\Rightarrow B$ 中，分别说明必要条件与充分条件，并解释为什么通过一个必要检查不能推出研究结论可信。
  \item 什么是合法反例？为什么 $w=(2,0,0)$ 不是“删除权重非负性后凸组合界失效”的受控反例？
\end{enumerate}

\subsection*{笔试推导}
\begin{enumerate}[nosep,leftmargin=*]
  \item 写出全称命题 $\forall x\in\mathcal D,A(x)\Rightarrow B(x)$ 的否定，并据此说明一个反例必须满足哪两件事。
  \item 不看正文，证明：若 $w_i\geq0$ 且 $\sum_iw_i=1$，则 $\min_i r_i\leq\sum_iw_ir_i\leq\max_i r_i$；同时给出上下界取等条件。
  \item 在 $\sum_iw_i=1$ 下证明：若凸组合界对所有 $r\in\mathbb R^n$ 成立，则每个 $w_i\geq0$。指出证明中哪里使用了“对所有”。
\end{enumerate}

\subsection*{数值编程}
\begin{enumerate}[leftmargin=*]
  \item 计算 $r=(0.02,-0.01,0.03)$、$w=(0.5,0.3,0.2)$ 的逐资产贡献，再把权重改为 $(1.5,-0.5,0)$。
  \item 画出 $2w-1$、$1-2w$ 及二者最小值，求 $0\leq w\leq1$ 上最坏收益的最大值。
  \item 两期方差都为 4，分别取协方差 $1,0,-1$，比较和的标准差。
\end{enumerate}

\subsection*{研究判断}
\begin{enumerate}[nosep,leftmargin=*]
  \item 研究员说“回测包含手续费，因此结果可信”。把这句话写成逻辑箭头，判断手续费检查是必要、充分还是充要条件，并列出两个仍可能失败的原因。
  \item 报告把 100 美元 P\&L、2\% 日收益和 15\% 年化波动率相加得到“综合得分”。逐项审查量纲和时间尺度，给出一个有意义的替代报告结构。
  \item 某人用一百万组随机收益都没有找到凸组合界反例，于是声称定理已被程序证明。说明有限实验为何不足以证明全称命题；给出一般证明，再构造删除非负权重条件后的反例。
\end{enumerate}
\endgroup

\subsection*{基础衔接：独立计算}
两期收益方差都为 $4$，协方差为 $1$。求收益之和的方差和标准差；若误用独立缩放，得到多少？


\subsection*{分级提示}
\begingroup\small
\begin{enumerate}[nosep,leftmargin=*]
  \item \textbf{口述概念：}先问一句话是否具有真假值；画出 $A\Rightarrow B$ 的箭头；合法反例要保留原命题全部假设，只让结论失败。检查 $w=(2,0,0)$ 的权重和。
  \item \textbf{笔试推导：}否定全称量词会变成存在量词；充分性分别展开 $R_p-m$ 与 $M-R_p$；必要性假设 $w_j<0$，再利用“对所有 $r$”选择 $r_j=0$、其余为 1。
  \item \textbf{数值编程：}先写出本章公式中的目标量，再显示中间计算。改变参数前预测结果方向，运行后说明差异来自模型、抽样还是数值近似。
  \item \textbf{研究判断：}手续费检查只能排除“完全忽略成本”；货币、无量纲收益与不同时间尺度的波动率不能直接相加；随机搜索适合寻找反例，不覆盖连续定义域。
\end{enumerate}
\endgroup


完整分层答案见配套答案册第 1 章。

\noindent\textbf{本章要点。}可证明的研究语言由对象、定义域、量词、假设和结论组成。必要与充分的方向不能混淆；证明方法要服从命题结构；反例必须隔离被删除的假设；量纲和时间尺度先于数值运算。掌握这些习惯，才有可能读懂后面更抽象的定义并判断它们何时适用于量化研究。
